\section{Dataset}

Two different datasets were used in this work: one dataset for model training
and fine-tuning, and a separate dataset for evaluation.

\subsection{Training Dataset}

The fine-tuned models were trained using the XLCoST dataset. XLCoST
(Cross-Lingual Code Snippet Dataset) is a large-scale benchmark designed for
cross-lingual code intelligence tasks, including code translation, code
summarization, and code search \cite{zhu2022xlcost}.

XLCoST contains parallel code snippets in seven programming languages:
C++, Java, Python, C\#, JavaScript, PHP, and C. In addition, the dataset also
includes aligned English descriptions. The dataset provides both snippet-level
and program-level parallel examples, making it suitable for supervised code
translation tasks. Given a code sample in one language, the dataset often
contains the same implementation in several other programming languages
\cite{zhu2022xlcost}.

For this thesis, only the Java-to-C++ portion of XLCoST was used. The training
notebooks load the original dataset splits and filter them to retain only Java
source code and the corresponding C++ target code. These code pairs are then
formatted into instruction-style prompts and used for supervised fine-tuning.

The Java-to-C++ translation direction was selected because both languages are
widely used, differ significantly in syntax and memory management, and are
commonly included in code translation benchmarks. In addition, Java and C++
examples are well represented in XLCoST, allowing the fine-tuned models to
learn from a sufficiently large number of aligned training examples.

One of the main advantages of XLCoST is its multilingual nature. Unlike smaller
translation datasets that only support one or two language pairs, XLCoST
contains fine-grained parallel data across multiple languages and supports ten
different cross-lingual tasks. The dataset contains more than one million
parallel code pairs across all supported languages, making it one of the
largest publicly available datasets for cross-lingual code intelligence
\cite{zhu2022xlcost}. 

\subsection{Evaluation Dataset}

A separate evaluation dataset was used to assess model performance after
training. Instead of evaluating only on the original XLCoST test split, this
thesis uses the CodeEditorBench framework for execution-based evaluation.

CodeEditorBench is a benchmark designed to evaluate code editing capabilities
of large language models in realistic software engineering scenarios
\cite{codeeditorbench}. It includes tasks such as debugging, code polishing,
requirement switching, and code translation. Unlike traditional benchmarks that
focus only on token-level similarity, CodeEditorBench evaluates whether the
generated code behaves correctly when executed.

For the purposes of this thesis, only the code translation subset of
CodeEditorBench was used. The models receive Java input code and are required
to generate the corresponding C++ translation. The generated output is then
compiled and tested using the execution framework provided by
CodeEditorBench.

The evaluation process was performed inside a Docker container using the
official CodeEditorBench environment. This ensures that all generated code is
tested under consistent conditions, including the same compiler versions,
runtime dependencies, and execution settings.

In addition to execution-based evaluation, BLEU score was also computed for the
generated translations. This makes it possible to compare similarity-based and
execution-based evaluation methods on the same set of examples.

Using separate datasets for training and evaluation helps reduce the risk of
data leakage and provides a more realistic assessment of model performance.
While XLCoST is primarily used to teach the model how to translate code,
CodeEditorBench is used to measure whether the translated code is actually
correct and executable.